\documentclass[12pt]{article}
\usepackage[margin=1in]{geometry}
\usepackage{amsmath, amssymb, amsfonts, bm}
\usepackage{enumitem}
\setlist[enumerate]{nosep}
\setlength{\parskip}{0.7em}
\setlength{\parindent}{0pt}

\begin{document}

\begin{center}
{\Large \textbf{Midterm Exam 1}} \\[0.9em]
\textbf{ECON 320 -- Introduction to Mathematical Economics} \\[0.7em]

October 01, 2025
 \\[0.8em]
Instructor: Muhebullah Karimzada
\end{center}

\hrule

\textbf{Instructions:}
\begin{itemize}
\item You are allowed \textbf{75 minutes} to complete the exam.
    \item Show all work. Unsupported answers receive little credit.

    \item Unless stated otherwise, all variables are real and differentiable as needed.
\end{itemize}

\hrule

\begin{enumerate}

\vspace{.7cm}
\item\textbf{ (7 points)} Let $U=\mathbb{R}$. Define
\[
A=\{x\in\mathbb{R}: x^2-7x+12\ge 0\}, \quad
B=\{x\in\mathbb{R}: |x-4|\le 3\}, \quad
C=\{x\in\mathbb{R}: 1<x<5\}.
\]

\begin{enumerate}[label=(\alph*)]
\item Write $A$, $B$, and $C$ explicitly as unions of intervals. \hfill (3)
\item Compute $A\cap B$, $A\cup C$, and $B\setminus C$. \hfill (2)
\item Decide whether $(A\cup B)\setminus C=(A\setminus C)\cup(B\setminus C)$. \hfill (2)
\end{enumerate}

\vspace{1cm}

\item \textbf{(8 points)} Let
\[
u(x,y)=x^2y - y^3 + \ln(1+xy).
\]
Suppose
\[
x = r^2,\quad y = s^3,\quad r>0,\ s>0.
\]

\begin{enumerate}[label=(\alph*)]
\item Compute $\partial u/\partial x$ and $\partial u/\partial y$. \hfill (3)
\item Using the chain rule, compute $\partial u/\partial r$ and $\partial u/\partial s$ in terms of $r,s$. \hfill (5)
\end{enumerate}
\vspace{1cm}
 \item \textbf{(8 points)} Let $F(x,y)=\ln(y+2x)-x^2+3y-10=0$ define $y$ as a function of $x$ near $x=1$.

\begin{enumerate}[label=(\alph*)]
\item Find $y$ such that $F(1,y)=0$ and verify the implicit function conditions at that point. \hfill (3)
\item Compute $\dfrac{dy}{dx}$ at that point. \hfill (3)
\item Approximate the change in $y$ when $x$ increases from $1$ to $1.04$ using the total differential. \hfill (2)
\end{enumerate}
\newpage
\item \textbf{(10 points)} Consider the transformation $(u,v,w)\mapsto(x,y,z)$:
\[
x = u^2 + 2vw,\quad
y = v^2 + 2uw,\quad
z = w^2 + 2uv.
\]

\begin{enumerate}[label=(\alph*)]
\item Compute the Jacobian matrix $J=\partial(x,y, z)/\partial(u,v,w)$. \hfill (5)
\item Evaluate $\det(J)$ at $(u,v,w)=(1,-2,3)$. Is the mapping locally invertible there? \hfill (5)
\end{enumerate}

\vspace{1cm}
\item\textbf{ (17 points)} A consumer maximizes utility
\[
U(x_1,x_2,\ell) = \ln(x_1) + \ln(x_2) + \beta \ln(\ell),
\]
where $x_1$ and $x_2$ are quantities of two goods and $\ell$ is leisure.  
The consumer faces the budget/time constraint:
\[
p_1 x_1 + p_2 x_2 + w\ell = I,
\]
where $p_1,p_2$ are prices, $w$ is the wage (price of leisure), and $I$ is income.

The first-order conditions are:
\[
\begin{cases}
F_1 = p_1 x_1 + p_2 x_2 + w\ell - I = 0, \\
F_2 = \dfrac{1}{x_1} - \lambda p_1 = 0, \\
F_3 = \dfrac{1}{x_2} - \lambda p_2 = 0.
\end{cases}
\]

\begin{enumerate}[label=(\alph*)]
\item Write the Jacobian $J=\partial(F_1,F_2,F_3)/\partial(x_1,x_2,\ell)$. \hfill (5)
\item Suppose $p_1=2$, $p_2=1$, $w=1$, $\beta=1$, and $I=12$. At the bundle $(x_1,x_2,\ell)=(2,4,6)$, compute $J$ and $\det(J)$. \hfill (7)
\item Using Cramer's Rule logic, determine the signs of $\dfrac{\partial x_1}{\partial I}$, $\dfrac{\partial x_2}{\partial I}$, $\dfrac{\partial \ell}{\partial I}$. Provide an economic interpretation. \hfill (5)
\end{enumerate}

\end{enumerate}

\end{document}
